\documentclass{article}
\usepackage[utf8]{inputenc}
\usepackage{geometry}
\geometry{a4paper, margin=1in}

\title{Can Unstable Relationships Become Sustainable?}
\author{Albert Jan van Hoek}
\date{\today}

\begin{document}

\maketitle

In my previous post, I suggested self-awareness as a possible way to resolve co-dependent and unequal relationships.

In response, I received the following question:

\begin{quote}
``I wonder how repairable these relationships are after an imbalance has occurred, and how visible the deep relationships and dependencies really are for others — for all players who also influence the structure.''
\end{quote}

This is a very good question. For me, it may be the central question of our time.

Tensions are rising — between people, within societies, and in our relationship with the planet. The real issue is whether unstable relationships can become sustainable again without escalating into conflict. Can we keep our peace.

\section*{Learning Networks and Survival}

I believe repair depends on whether we as individuals, and our systems can recognize themselves as learning networks.

A network survives not by remaining unchanged, but by adapting.
A brain stays intact because it continuously adjusts.

But a brain can only survive if the conditions it depends on also survive.

Without food, there is no metabolism.
Without metabolism, no brain.

The same holds for relationships and societies.

They persist only if the ecological and social foundations they depend on are preserved.

\section*{Learning as Protection}

Learning, therefore, is not just internal adjustment.
It is the capacity to recognize and protect one’s own dependencies.

Long-term stability emerges when a system adapts in ways that keep the network — and the conditions that sustain it — intact.

\section*{Peace as a Condition for Learning}

Peace, in this sense, is not merely the absence of conflict.

It is the condition under which interdependent systems can continue to learn — and therefore continue to exist.

This is what I call the \textbf{Theory of Long-Term Collaboration}.

And how visible those dependencies are?
That depends on whether we choose to make them visible.

\end{document}
